
\subsection{$A_\infty$ 双模同态的 Whitehead 定理}
\label{subsec:5-2-6}

本小节借助数目滤过\liftIndex{I094}，
证明定理 5.2.35 的非滤过版本
\liftIndex{I192}。

令 $\overline C_i,\overline C'_i$
（$i=0,1$）为 $A_\infty$ 代数，
$\overline D$ 与 $\overline D'$
分别为
$(\overline C_1,\overline C_0)$
与
$(\overline C'_1,\overline C'_0)$
上的双模。
设
$\overline{\mathfrak f}^{(i)}:
\overline C_i\to\overline C'_i$
是同伦等价，
其同伦逆为
$(\overline{\mathfrak f}^{(i)})^{-1}$
（$i=0,1$）。
假设位于
$(\overline{\mathfrak f}^{(1)},
  \overline{\mathfrak f}^{(0)})$
之上的 $A_\infty$ 双模同态
$\overline\varphi:
\overline D\to\overline D'$
在 $\overline{\mathfrak n}_{0,0}$
上同调上诱导同构。
我们将证明存在位于
$((\overline{\mathfrak f}^{(1)})^{-1},
   (\overline{\mathfrak f}^{(0)})^{-1})$
之上的同伦逆
$\overline\varphi':
\overline D'\to\overline D$。

为此，将对 $K,L$ 作归纳，
证明下列命题。
令
$\overline{\mathfrak h}^{(i)}:
\overline C_i\to
\overline{\mathfrak C}_i$
是从恒等同态到
$(\overline{\mathfrak f}^{(i)})^{-1}
\circ\overline{\mathfrak f}^{(i)}$
的同伦，
并令 $\overline{\mathfrak D}$
为 $[0,1]\times\overline D$ 的模型。
注意，$\overline{\mathfrak D}$
与 $\overline{\mathfrak h}^{(i)}$
均可任意选取；
这一点对于证明定理 5.2.35
最后一项结论很重要。

\begin{proposition}{5.2.48}\label{prop:5-2-48}
对任意 $K,L=1,2,\ldots$，
均存在位于
$((\overline{\mathfrak f}^{(1)})^{-1},
   (\overline{\mathfrak f}^{(0)})^{-1})$
之上的 $A_{K,L}$ 双模同态
$\overline\varphi'_{(K,L)}:
\overline D'\to\overline D$，
以及位于
$(\overline{\mathfrak h}^{(1)},
  \overline{\mathfrak h}^{(0)})$
之上的 $A_{K,L}$ 双模同态
$\overline{\mathfrak H}_{(K,L)}:
\overline D\to\overline{\mathfrak D}$，
使得：
\begin{enumerate}[label=\textup{(5.2.49.\arabic*)},leftmargin=5.8em]
\item\label{cond:5-2-49-1}
若 $K'\leq K$ 且 $L'\leq L$，
则
$\widehat{\overline\varphi}'_{(K,L)}$
在
\[
B_{1\cdots K'}\overline C'_1[1]
\otimes\overline D'[1]
\otimes B_{1\cdots L'}\overline C'_0[1]
\]
上与
$\widehat{\overline\varphi}'_{(K',L')}$
相同；而
$\widehat{\overline{\mathfrak H}}_{(K,L)}$
在
\[
B_{1\cdots K'}\overline C_1[1]
\otimes\overline D[1]
\otimes B_{1\cdots L'}\overline C_0[1]
\]
上与
$\widehat{\overline{\mathfrak H}}_{(K',L')}$
相同。
\item\label{cond:5-2-49-2}
\[
\mathrm{Eval}_{s=0}
\circ\overline{\mathfrak H}
=\mathrm{id},
\qquad
\mathrm{Eval}_{s=1}
\circ\overline{\mathfrak H}
=\widehat{\overline\varphi}'_{(K,L)}
\circ\widehat{\overline\varphi}.
\]
\end{enumerate}
\end{proposition}

\begin{proof}
对 $K,L$ 归纳。
假设命题在 $(K,L)$ 成立，
要证明 $(K+1,L)$ 与
$(K,L+1)$ 两种情形；
只写前者。

考虑
$\overline{\mathfrak H}_{(K,L)}:
\overline D\to\overline{\mathfrak D}$。
由于
$\mathrm{Eval}_{s=0}
\circ\overline{\mathfrak H}_{(K,L)}$
是恒等同态，且
$\mathrm{Eval}_{s=0}$
在上同调上诱导同构，
由引理 5.2.46，
可找到
$\overline{\mathfrak H}^{0}_{K+1,L}$，
将
$\overline{\mathfrak H}_{(K,L)}$
延拓为
$\overline{\mathfrak H}^{0}_{(K+1,L)}$。
还可假设
\[
\mathrm{Eval}_{s=0}
\circ\overline{\mathfrak H}^{0}_{(K+1,L)}
=\mathrm{id}.
\]
这一点可按命题 4.5.6
证明开头的类似论证得到。

另一方面，
$\widehat{\overline\varphi}'_{(K,L)}
\circ\widehat{\overline\varphi}$
同伦于恒等同态，
而 $\widehat{\overline\varphi}$
在上同调上诱导同构。
由引理 5.2.46，
存在延拓
$\overline\varphi'_{(K,L)}$
的
$\overline\varphi'_{(K+1,L)}$。

考虑差
\translatornote{原书实体页 302 在定义
$\Xi$ 时把输入张量统写为
$\overline\varphi_1\otimes\cdots
\otimes\overline\varphi_1$；
但紧接着明确写出左右代数输入
分别取
$(\overline{\mathfrak f}^{(1)}_1)^{\otimes(K+1)}$
与
$(\overline{\mathfrak f}^{(0)}_1)^{\otimes L}$，
中间双模输入才取
$\overline\varphi_1$。
此处依后式展开；见疑误 A-0286。}
\[
\begin{aligned}
\Xi={}&-
\mathrm{Eval}_{s=1}\circ
\overline{\mathfrak H}^{0}_{(K+1,L)}
\\
&+\widehat{\overline\varphi}'_{(K+1,L)}
\circ\bigl(
(\overline{\mathfrak f}^{(1)}_1)^{\otimes(K+1)}
\otimes\overline\varphi_1
\otimes
(\overline{\mathfrak f}^{(0)}_1)^{\otimes L}
\bigr).
\end{aligned}
\]
$\Xi$ 是
\[
\operatorname{Hom}\bigl(
B_{K+1}\overline C_1[1]
\otimes\overline D[1]
\otimes B_L\overline C_0[1],
\overline D[1]\bigr)
\]
中的 $\delta_1$-上闭链。
由于
\[
(\overline{\mathfrak f}^{(1)}_1)^{\otimes(K+1)}
\otimes\overline\varphi_1
\otimes
(\overline{\mathfrak f}^{(0)}_1)^{\otimes L}
\]
是上链同伦等价，
存在 $\delta_1$-上闭链
\[
\Delta\overline\varphi'_{(K+1,L)}
\in\operatorname{Hom}\bigl(
B_{K+1}\overline C'_1[1]
\otimes\overline D'[1]
\otimes B_L\overline C'_0[1],
\overline D[1]\bigr),
\]
使得
\[
\Bigl[
\Xi+
\Delta\overline\varphi'_{(K+1,L)}
\circ
\bigl(
(\overline{\mathfrak f}^{(1)}_1)^{\otimes(K+1)}
\otimes\overline\varphi_1
\otimes
(\overline{\mathfrak f}^{(0)}_1)^{\otimes L}
\bigr)
\Bigr]=0.
\]
以
$\overline\varphi'_{(K+1,L)}
+\Delta\overline\varphi'_{(K+1,L)}$
替换原来的
$\overline\varphi'_{(K+1,L)}$，
即可假设 $[\Xi]=0$。
因此，存在
\[
\Delta_1\overline{\mathfrak H}_{(K+1,L)}
\in\operatorname{Hom}\bigl(
B_{K+1}\overline C_1[1]
\otimes\overline D[1]
\otimes B_L\overline C_0[1],
\overline D[1]\bigr),
\]
满足
\[
\begin{aligned}
&\delta_1\bigl(
\Delta_1\overline{\mathfrak H}_{(K+1,L)}
\bigr)
\\
&\quad=
\overline\varphi'_{(K+1,L)}
\circ\bigl(
(\overline{\mathfrak f}^{(1)}_1)^{\otimes(K+1)}
\otimes\overline\varphi_1
\otimes
(\overline{\mathfrak f}^{(0)}_1)^{\otimes L}
\bigr)
-
(\mathrm{Eval}_{s=1})_1
\circ\overline{\mathfrak H}^{0}_{(K+1,L)}.
\end{aligned}
\]
由条件 (5.2.22.3)，
可找到
\[
\Delta\overline{\mathfrak H}_{(K+1,L)}
\in\operatorname{Hom}\bigl(
B_{K+1}\overline C_1[1]
\otimes\overline D[1]
\otimes B_L\overline C_0[1],
\overline{\mathfrak D}[1]\bigr),
\]
使得
\[
\begin{aligned}
(\mathrm{Eval}_{s=0})_1
\circ\Delta\overline{\mathfrak H}_{(K+1,L)}
&=0,\\
(\mathrm{Eval}_{s=1})_1
\circ\Delta\overline{\mathfrak H}_{(K+1,L)}
&=\Delta_1
\overline{\mathfrak H}_{(K+1,L)}.
\end{aligned}
\]
此时直接检查可知
\[
\overline{\mathfrak H}_{(K+1,L)}
=\overline{\mathfrak H}^{0}_{(K+1,L)}
+\delta_1\bigl(
\Delta\overline{\mathfrak H}_{(K+1,L)}
\bigr)
\]
具有所需性质。
\end{proof}

\begin{proof}[定理 5.2.35 的非滤过版本]
如 $A_\infty$ 代数情形一样，
由命题 5.2.48 即得结论。
\end{proof}

在进入滤过版本之前，
再说明定理 5.2.35 的
非滤过含幺版本如何证明。
主要区别在于使用
\[
\overline C_1[1]\otimes
B_{K-1}\overline C_{1,\mathrm{red}}[1]
\otimes\overline D[1]\otimes
B_{L-1}\overline C_{0,\mathrm{red}}[1]
\otimes\overline C_0[1].
\]
这里对幺元 $e$ 记
$\overline C_{\mathrm{red}}
=\overline C/R[e]$。
于是可按第 4.4.4 小节的类似论证
修改定理 5.2.35 的证明；
细节留给读者。
